\section{Introducción}

El presente documento describe la propuesta de desarrollo de \textbf{SocioUnido}, una plataforma inteligente orientada a la gestión administrativa, fidelización y modernización operativa de clubes de fútbol argentinos. El objetivo principal del proyecto es brindar una solución tecnológica integral que permita optimizar la administración de socios y espacios rentables, reducir pérdidas económicas derivadas de la morosidad y el fraude en accesos, y mejorar la interacción entre las instituciones deportivas y sus hinchas.

Actualmente, una gran cantidad de clubes de fútbol, especialmente aquellos pertenecientes al ascenso argentino y categorías intermedias, continúan dependiendo de procesos manuales, sistemas fragmentados o herramientas administrativas desactualizadas. Esta situación genera ineficiencias operativas, dificultades para escalar servicios digitales, pérdida de oportunidades de monetización y una experiencia deficiente para el socio. Asimismo, los días de partido representan un escenario crítico debido a problemas de conectividad, validación de accesos y sobrecarga administrativa.

Frente a este contexto, SocioUnido propone una arquitectura tecnológica moderna basada en un modelo de \textit{software como servicio (Software as a Service - SaaS)}, complementada con herramientas de inteligencia artificial, automatización conversacional y mecanismos de validación fuera de línea (offline) mediante códigos dinámicos. La plataforma busca transformar el enfoque tradicional de gestión reactiva en un modelo predictivo y proactivo, permitiendo a los clubes anticipar situaciones de morosidad, automatizar procesos administrativos y mejorar la seguridad operativa en eventos masivos.

El presente anteproyecto expone de manera estructurada los distintos aspectos conceptuales, técnicos y metodológicos involucrados en el desarrollo de la solución. En primer lugar, se presenta un relevamiento del estado actual del mercado y de las tecnologías utilizadas en sistemas similares, identificando las principales limitaciones existentes en las soluciones actualmente adoptadas por las instituciones deportivas. Posteriormente, se desarrolla el problema detectado y se fundamenta la necesidad de una herramienta específicamente diseñada para la realidad operativa y económica del fútbol argentino.

A continuación, se describe la solución propuesta, incluyendo sus principales módulos funcionales, los componentes tecnológicos involucrados y las decisiones de diseño preliminares consideradas para garantizar escalabilidad, seguridad y capacidad de integración. Además, se incorpora una evaluación preliminar del impacto económico, social y operativo que podría generar la implementación de la plataforma tanto para los clubes como para los socios usuarios del sistema.

Finalmente, el documento detalla la metodología de trabajo seleccionada para el desarrollo del proyecto, las estrategias de validación y control de calidad previstas, y el plan de actividades estimado para las distintas etapas de construcción, prueba e implementación del producto. De esta manera, se busca establecer una base sólida para el desarrollo de un sistema innovador, técnicamente robusto y alineado con las necesidades reales del ecosistema deportivo argentino.
